
% BEGIN bib
 
\par\medskip\noindent\textbf{References}\par
  
\par\noindent  [1]  Radhakisan Baheti and Helen Gill.  Cyber-physical systems.  The impact of control technology, 12(1):161–166, 2011.  
 
\par\noindent  [2]  Nicolas Behr and Jean Krivine.  Rewriting theory for the life sciences.  In International Conference on Graph Transformation, pages 185–202, 2020.  
 
\par\noindent  [3]  Jasmin Christian Blanchette, Andrei Popescu, and Dmitriy Traytel.  Cardinals in Isabelle/HOL.  In Lecture Notes in Computer Science, volume 8558, pages 111–127, 2014.  doi:10.1007/978-3-319-08970-6\_8.  
 
\par\noindent  [4]  Mirna Bognar.  A point version of decreasing diagrams.  In Proceedings Accolade 1996. Dutch Graduate School in Logic, pages 1–14, 1997.  
 
\par\noindent  [5]  F. Bonchi, F. Gadducci, A. Kissinger, P. Sobociński, and F. Zanasi.  String diagram rewrite theory I: Rewriting with Frobenius structure.  Journal of the ACM, 69(2):14:1–14:58, 2022.  doi:10.1145/3502719.  
 
\par\noindent  [6]  N.G. de Bruijn.  A note on weak diamond properties.  Memorandum 78–08, Eindhoven Uninversity of Technology, 1978.  
 
\par\noindent  [7]  R. Diestel and O. Pikhurko.  On the cofinality of infinite partially ordered sets: Factoring a poset into lean essential subsets.  Order, 20:53–66, 2003.  doi:10.1023/A:1024449306316.  
 
\par\noindent  [8]  Jörg Endrullis and Jan Willem Klop.  De Bruijn’s weak diamond property revisited.  Indagationes Mathematicae, 24:1050–1072, 2013.  
 
\par\noindent  [9]  Jörg Endrullis, Jan Willem Klop, and Roy Overbeek.  Decreasing diagrams with two labels are complete for confluence of countable systems.  In 3rd International Conference on Formal Structures for Computation and Deduction (FSCD 2018), pages 14:1–14:15. Schloss Dagstuhl – Leibniz-Zentrum für Informatik, 2018.  doi:10.4230/LIPICS.FSCD.2018.14.  
 
\par\noindent  [10]  Jörg Endrullis, Jan Willem Klop, and Roy Overbeek.  Decreasing diagrams for confluence and commutation.  Logical Methods in Computer Science, 16:23:1–23:25, 2020.  doi:10.23638/LMCS-16(1:23)2020.  
 
\par\noindent  [11]  A.A. Fraenkel, Y. Bar-Hillel, and A. Levy.  Foundations of Set Theory.  Elsevier, 1973.  
 
\par\noindent  [12]  Jean Goubault-Larrecq.  Non-Hausdorff topology and domain theory: Selected topics in point-set topology.  Cambridge University Press, 2013.  
 
\par\noindent  [13]  N. Hirokawa, A. Middeldorp, C. Sternagel, and S. Winkler.  Abstract completion, formalized.  Logical Methods in Computer Science, 15(3):19:1–19:42, 2019.  doi:10.23638/LMCS-15(3:19)2019.  
 
\par\noindent  [14]  Isabelle proof assistant.  http://isabelle.in.tum.de.  [Online].  
 
\par\noindent  [15]  J.R. Isbell.  The category of cofinal types. II.  Trans. Amer. Math. Soc., 116:394–416, 1965.  
 
\par\noindent  [16]  Ievgen Ivanov.  Generalized Newman’s lemma for discrete and continuous systems.  In Leibniz International Proceedings in Informatics, volume 260, pages 9:1–9:17, 2023.  doi:10.4230/LIPICS.FSCD.2023.9.  
 
\par\noindent  [17]  Ievgen Ivanov.  Formal proof of completeness of the decreasing diagrams method for proving confluence of relations of the least uncountable cardinality (supplementary material for this paper).  https://doi.org/10.5281/zenodo.14254256, 2024.  [Online].  
 
\par\noindent  [18]  Ievgen Ivanov.  On Newman’s lemma and non-termination.  In CEUR-WS.org, volume 3624, pages 14–24, 2024.  URL: https://ceur-ws.org/Vol-3624/Paper\_2.pdf.  
 
\par\noindent  [19]  Ievgen Ivanov.  On non-triviality of the hierarchy of decreasing Church-Rosser abstract rewriting systems.  In Proceedings of the 13th International Workshop on Confluence, pages 30–35, 2024.  
 
\par\noindent  [20]  Ievgen Ivanov.  On the cofinality property in the context of the hierarchy of decreasing Church-Rosser abstract rewriting systems.  In Information Technology and Implementation, November 20–21, 2024, Kyiv, Ukraine, 2024.  
 
\par\noindent  [21]  Ievgen Ivanov.  Completeness of decreasing diagrams for the least uncountable cardinality.  Archive of Formal Proofs, April 2025.   https://isa-afp.org/entries/Completeness\_Decreasing\_Diagrams\_for\_N1.html, Formal proof development.  
 
\par\noindent  [22]  Ievgen Ivanov.  Formalization of a confluent abstract rewriting system of the least uncountable cardinality outside of the class DCR2.  https://doi.org/10.5281/zenodo.14740062, 2025.  [Online].  
 
\par\noindent  [23]  John L. Kelley.  General topology.  Van Nostrand, 1955.  
 
\par\noindent  [24]  Jan Willem Klop.  Combinatory Reduction Systems.  PhD thesis, Rijks-universiteit Utrecht, 1980.  
 
\par\noindent  [25]  Alexander Krauss and Andreas Schropp.  A mechanized translation from Higher-Order Logic to set theory.  In Lecture Notes in Computer Science, volume 6172, pages 323–338. Springer, 2010.  doi:10.1007/978-3-642-14052-5\_23.  
 
\par\noindent  [26]  Ondřej Kunčar and Andrei Popescu.  Comprehending Isabelle/HOL’s consistency.  In Lecture Notes in Computer Science, volume 10201, pages 724–749. Springer, 2017.  doi:10.1007/978-3-662-54434-1\_27.  
 
\par\noindent  [27]  Ondřej Kunčar and Andrei Popescu.  A consistent foundation for Isabelle/HOL.  Journal of Automated Reasoning, 62:531–555, 2019.  doi:10.1007/S10817-018-9454-8.  
 
\par\noindent  [28]  Edward A. Lee.  Fundamental limits of cyber-physical systems modeling.  ACM Transactions on Cyber-Physical Systems, 1(1):1–26, 2016.  doi:10.1145/2912149.  
 
\par\noindent  [29]  George Markowsky.  Chain-complete posets and directed sets with applications.  Algebra universalis, 6(1):53–68, 1976.  
 
\par\noindent  [30]  Elliott Mendelson.  Introduction to Mathematical Logic.  CRC Press, 2015.  
 
\par\noindent  [31]  Eric C. Milner and Maurice Pouzet.  On the cofinality of partially ordered sets.  Ordered sets, pages 279–298, 1982.  
 
\par\noindent  [32]  Eric C. Milner and Karel Prikry.  The cofinality of a partially ordered set.  Proc. London Math. Soc., 46:454–470, 1983.  
 
\par\noindent  [33]  Chuangjie Xu Nicolai Kraus, Fredrik Nordvall Forsberg.  Type-theoretic approaches to ordinals.  Theoretical computer science, 957, 2023.  doi:10.1016/J.TCS.2023.113843.  
 
\par\noindent  [34]  Tobias Nipkow, Markus Wenzel, and Lawrence C Paulson.  Isabelle/HOL: a proof assistant for higher-order logic.  Springer, 2002.  doi:10.1007/3-540-45949-9.  
 
\par\noindent  [35]  Andrew Pitts.  The HOL logic.  In Introduction to HOL: A Theorem Proving Environment for Higher Order Logic, pages 191–232. Cambridge University Press, 1993.  
 
\par\noindent  [36]  Raymond Smullyan.  Set theory and the continuum problem.  Clarendon Press, 1996.  
 
\par\noindent  [37]  Christian Sternagel and Rene Thiemann.  Abstract rewriting.  Archive of Formal Proofs, June 2010.  
 
\par\noindent  [38]  Stevo Todorcevic.  Directed sets and cofinal types.  Trans. Amer. Math. Soc., 290:711–723, 1985.  
 
\par\noindent  [39]  John W. Tukey.  Convergence and uniformity in topology.  Princeton University Press, 1940.  
 
\par\noindent  [40]  Vincent Van Oostrom.  Confluence by decreasing diagrams.  Theoretical computer science, 126(2):259–280, 1994.  doi:10.1016/0304-3975(92)00023-K.  
 
\par\noindent  [41]  Vincent Van Oostrom.  Confluence for Abstract and Higher-Order Rewriting.  PhD thesis, Vrije Universiteit Amsterdam, 1994.  
 
\par\noindent  [42]  Vincent Van Oostrom.  Confluence by decreasing diagrams – converted.  In Lecture Notes in Computer Science, volume 5117, pages 306–320, 2008.  
 
\par\noindent  [43]  Shang-Zhi Wang and Bo Yu Li.  On the minimal cofinal subsets of a directed quasi-ordered set.  Discrete Mathematics, 48:289–306, 1984.  doi:10.1016/0012-365X(84)90189-4.  
 
\par\noindent  [44]  Freek Wiedijk.  The seventeen provers of the world: Foreword by Dana S. Scott, volume 3600.  Springer, 2006.  doi:10.1007/11542384.  
 
\par\noindent  [45]  Harald Zankl.  Confluence by decreasing diagrams – formalized.  In 24th International Conference on Rewriting Techniques and Applications, pages 352–367. Leibniz International Proceedings in Informatics, 2013.  doi:10.4230/LIPICS.RTA.2013.352.  
 
\par\noindent  [46]  Harald Zankl.  Decreasing diagrams.  Archive of Formal Proofs, 2013.  
  
% END bib
